\scanpage{17}
On a un foncteur
\[
C'_0\longrightarrow C''_0
\]
C'est une équivalence de catégories \uuline{si $n\geq2$}

(mais \uline{non si $n=2$})

En effet, loc$^{\mathrm t}$ on peut écrire $C=\check{\mathbb P}(V)$, et remonter de façon \uline{unique} les \inserted{ps. réflex.} $\tau_i^C$ en des \inserted{ps. réfl.} $\tau_i^V$, qui vont satisfaire aux m$^{\mathrm s}$ conditions de commutation. \margin{\unsure{$\lambda$} (car $n\geq2$!)} Le module $V$ est déterminé modulo $\otimes L$, $L$ inv., et cette indétermination est levée par la rigidification $L_0\simeq\mathcal O_S$ \ldots

\deleted{\uline{NB}} \uline{Rescholie}

\deleted{\unsure{Donc}} Si $n\geq2$, on a \unsure{donc|dans}
\[
C_0\mathrel{\underset{\rule{1.4em}{.35pt}}{\approx}}C'_0\mathrel{\underset{\rule{1.4em}{.35pt}}{\approx}}C''_0
\]

Formes quadratiques \deleted{\unsure{pour}} invariantes sur $V$

\margin{\uuline{$n\geq2$}} \deleted{Si $n\geq$} Celles qui sont $\ne0$ sur tte fibre correspondent $1-1$ aux hyperquadriques \inserted{$Q_C$} de $C=\mathrm P(V)$, invariantes par $\varphi_{\unsure{\mathscr G|G|\mathfrak G}}^C$. On suppose ici \uline{$\lambda_i=-1$} $\forall i$ \margin{donc on remplace $\mathscr G$ par $G$, \unsure{avec|ou} les relations $\underline{\tau}_i^2=1$} % slanted
(condition néc. pour qu'il existe une $f$ \deleted{\unsure{non nulle}} inv. \unsure{avec|non|non dég.} $f_V(a_i)$ inv. \ldots) Il faut écrire
\[
\underset{\substack{\|\\c_i}}{\unsure{f_V|\varphi_V}(a_i)}\,a_i^\vee+\psi_V(a_i)=0\qquad 0\leq i\unsure{<|\leq}n
\]
i.e.
\[
\varphi_V(a_i,a_j)=-c_i\langle a_j,a_i^\vee\rangle\qquad 0\leq i,j\leq n
\]
donc
\[
\text{i)}\qquad\varphi_V(a_i,a_j)=0\qquad i,j\ \text{non consécutifs.}
\]
% archive sheet number: 16

